\NewDocumentCommand{\DrawEinsteinSolid}{ O{5} O{5} O{0.4} }{
\def\L{3}
\def\DeltaL{3.05}
\def\Lx{#1}
\def\Ly{#1}
\pgfmathsetmacro{\Temp}{#3}

\foreach \j in {1,...,\Lx}{
\foreach \i in {1,...,\Ly}{
    \begin{scope} [shift={(\i*\DeltaL,\j*\DeltaL)}]
        \pgfmathsetmacro{\DeltaX}{\Temp*rand}
        \DrawLatticeCell{0,0}[\L]
        \draw[spring, segment length=(1+\DeltaX)*2.0] (-\L/2+0.03,0) -- (-0.2+\DeltaX\L,0);
        \DrawParticle{\DeltaX\L,0}
    \end{scope}
}}

\begin{scope}[shift={(+\DeltaL/2, +\DeltaL/2)}]
        \draw[color=PrimaryColor, thick, step=\DeltaL] 
            (0, 0) grid (\Lx*\DeltaL, \Ly*\DeltaL);
\end{scope}
}

\NewDocumentCommand{\EinsteinSolid}{ O{5} O{5} O{0.4} }{
    \begin{tikzpicture}[x=\baselineskip, y=\baselineskip]
        \DrawEinsteinSolid[#1][#2][#3]
    \end{tikzpicture}
}

\newcommand{\OscilatorInAir}{
    \fill[pattern=north east lines] (-1,-1.25) rectangle (-1.25,1.25);
    \draw[line width=0.4] (-1,-1.25) -- (-1,1.25);
    \begin{scope}[shift={(-1,0)}]
        \pgfmathsetmacro{\DeltaX}{0.8*rand}
        \draw[spring, segment length=(1+\DeltaX)*2.0] (0,0) -- (\L/2-0.2+\DeltaX\L-0.03,0);
        \DrawParticle{\DeltaX\L-0.03+\L/2,0}
    \end{scope}
}

\newcommand{\EinsteinSolidMultiplicity}{
\begin{tikzpicture}[x=\baselineskip, y=\baselineskip, 
        trim left=-8\baselineskip, 
        trim right=8\baselineskip ]
\def\L{3}
\def\DeltaL{3.05}
\def\N{6}       % Quantidade de bolinhas na base da pirâmide
\def\dist{1.05} % Distância D entre os centros das bolinhas 
\def\MaxY{\dist*\N*0.8660}
\pgfmathsetmacro{\Radius}{0.05*\baselineskip}
\def\MinY{-\Radius}
\def\MidY{0.5*\MaxY+0.5*\MinY}
\def\MinXPos{-7.5}
\def\MaxXPos{+7.5}
\def\MinX{\MinXPos-\N*\Radius}

\begin{scope}[shift={(\MinXPos,0)}]
    \foreach \camada in {1, ..., \N} {
        \pgfmathsetmacro{\MaxIndex}{\N - \camada}
        \foreach \i in {0, ..., \MaxIndex} {
            \pgfmathsetmacro{\Y}{(\camada - 1) * \dist * 0.8660}
            \pgfmathsetmacro{\X}{(\i - \MaxIndex/2) * \dist}
            \DrawParticle[MediumGreen]{\X, \Y}[\baselineskip]
        }
    }
    \pgfmathsetmacro{\Y}{(\N-1)*\dist*0.8660}
    \pgfmathsetmacro{\X}{0*\dist}
    \DrawParticle[MediumRed]{\X, \Y}[\baselineskip]
    \pgfmathsetmacro{\Y}{\Y-1*\dist*0.8660}
    \pgfmathsetmacro{\X}{-0.5*\dist}
    \DrawParticle[MediumRed]{\X, \Y}[\baselineskip]
    \pgfmathsetmacro{\X}{0.5*\dist}
    \DrawParticle[MediumRed]{\X, \Y}[\baselineskip]
\end{scope}

\begin{scope}[shift={(0,\MidY)}]
    \fill[pattern=north east lines] (-1,-1.25) rectangle (-1.25,1.25);
    \draw[line width=0.4] (-1,-1.25) -- (-1,1.25);
    \begin{scope}[shift={(-1,0)}]
        \pgfmathsetmacro{\DeltaX}{0.8*rand}
        \draw[spring, segment length=(1+\DeltaX)*2.0] (0,0) -- (\L/2-0.2+\DeltaX\L-0.03,0);
        \DrawParticle{\DeltaX\L-0.03+\L/2,0}
    \end{scope}
\end{scope}

\def\L{0.9}
\def\DeltaL{0.905}
\def\MaxX{\MaxXPos+\DeltaL*5}
\def\IdxList{-2,-1,0,1,2}
\begin{scope}[shift={(\MaxXPos,\MidY)}]
    \foreach \j in \IdxList{
    \foreach \i in \IdxList{
        \begin{scope} [shift={(\i*\DeltaL,\j*\DeltaL)}]
            \DrawLatticeCell{0,0}[\L]
        \end{scope}
    }}
\begin{scope}[shift={(-2.5*\DeltaL, -2.5*\DeltaL)}]    
    \draw[color=DarkBlue, thick, step=\DeltaL]
            (0, 0) grid (5*\DeltaL,5*\DeltaL);
\end{scope}
\end{scope}

\draw [-{Triangle[length=0.5em, width=0.4em]}, very thick, color=black]
    (1.2*\MinX,1.2*\MaxY) to [bend right=-40] node[above=1em, align=center] {Escolhemos os\\ quanta de energia} (\MinXPos,\MaxY);

\pgfmathsetmacro{\NewY}{(\N-1.5)*\dist*0.8660}
\draw [-{Triangle[length=0.5em, width=0.4em]}, very thick, color=black]
    (0.5*\MinX+2.5,\NewY) to node[anchor=west, xshift=-3em, yshift=1.5em, align=center] {Dotamos um oscilador\\ com esta energia} (0.5*\MinX+5,\NewY);

\draw [-{Triangle[length=0.5em, width=0.4em]}, very thick, color=black]
    (+0.8,0.5*\MidY) to [bend right=40] node[below=1em, align=center] {Fixamos o oscilador em um\\ retículo vago da rede}(\MaxXPos-3,0.5*\MidY);

\end{tikzpicture}
}